\subsection{\problemblock{*When} Data Block}\label{sec:block-when}

Segment initial conditions are obtained from the trajectory initial conditions or from
the final conditions of the previous segment. Segments are flown according to the
guidance rules given in the \problemblock{*fly} or \problemblock{*rail} data
blocks and they continue until a final condition is met. At that point, the segment
ends, and if there is another segment in the trajectory, it is started.

The \problemblock{*when} data block is used to provide a segment final condition.
It also tells TAOS what to do next, that is, whether to transfer to another segment or
to end the trajectory. Each segment must have at least one
\problemblock{*when} data block giving a final condition; otherwise, the segment
never terminates.

An example of a segment final condition is

\begin{lstlisting}[style=taosmanual]
*when alt<200000 goto 4
\end{lstlisting}

which reads ``end the current segment when altitude becomes less than
\SI{200000}{\foot} and then start executing segment number 4.'' The variable name,
in this case \statevar{alt}, can be any valid output-variable name including
user-defined output variables. A list of output variables is given in the Appendix.

Final conditions can also compare two output variables, for example,

\begin{lstlisting}[style=taosmanual]
*when alpha=alpha_l/d goto 2
\end{lstlisting}

ends the segment when the angle of attack equals the angle of attack for maximum
lift-to-drag ratio.

\taossubhead{Values from Other Trajectories}

Output variables from other trajectories can be referenced by using subscripts, for
example,

\begin{lstlisting}[style=taosmanual]
*when mach[3]>4 goto 21
\end{lstlisting}

says to end the segment when the Mach number on trajectory 3 becomes greater
than 4.0. Subscripts are always enclosed in square brackets.

When final conditions are used to compare two variables, subscripts can be used on
both variables to compare values from different trajectories. For example,

\begin{lstlisting}[style=taosmanual]
*when east[2]=east[1] goto 7
\end{lstlisting}

ends the segment when the two trajectories have the same east position.

\taossubhead{Final Condition Relationships}

The relationship in a \problemblock{*when} data block can be greater than, less
than, or equal to. Since all values are continuous real numbers, using the greater-than
or less-than relationship often has the same effect as the equal sign because the
segment will end when the variable is equal to the value. The difference occurs at
the beginning of a segment.

When a segment starts executing, it checks the final conditions for immediate
termination. If any of the final conditions are satisfied initially, they cause an
immediate transfer to another segment. As a result the final condition
\texttt{mach>3} is much different from \texttt{mach=3}. A \texttt{mach>3} final
condition terminates the segment immediately if the Mach number is greater than
3.0, whereas the \texttt{mach=3} final condition only terminates the segment
immediately if Mach is exactly equal to 3.0. In practice use of greater-than and
less-than is more common than the equal sign.

\taossubhead{Ending Trajectories}

So far, all of the examples have shown segment transfers using the \texttt{goto}
keyword. When the end of a trajectory is reached, there is no segment to transfer to,
so the \texttt{stop} keyword is used as follows:

\begin{lstlisting}[style=taosmanual]
*when time=65 stop
\end{lstlisting}

This data block ends the trajectory when the time is equal to 65 seconds.
Calculations proceed until a segment final condition with \texttt{stop} is
encountered.

\taossubhead{Multiple Final Conditions}

Segments can have any number of final conditions. The first condition encountered,
while computing the trajectory, is executed. This allows complex branching in the
trajectory depending on what happens as shown in the following example:

\begin{lstlisting}[style=taosmanual]
*when tmark>25 goto 6
*when alphat>5 goto 11
\end{lstlisting}

If the ``mark'' time becomes equal to 25 seconds before the total angle of attack
reaches $5^{\circ}$, then segment 6 is executed next; otherwise, segment 11 is
executed next.

If two final conditions are immediately satisfied at the beginning of a segment, the
first one listed is executed.

It is good practice to have at least one final condition on a time variable. For example,

\begin{lstlisting}[style=taosmanual]
*when gamgd<0 goto 4
*when tseg>1000 goto 10
\end{lstlisting}

normally ends the segment when the flight-path angle becomes zero. But if there is
a problem and the flight-path angle never reaches zero, the segment still ends after
1000 seconds. This prevents infinite loops.
